\section{对格点计算的影响}
\label{sec:lattice}

我们在\sec{ope}中的证明澄清了重正化准 PDF、\spcorr 与赝 PDF 分布之间的关系。作为实际问题，把这些方程用于将\spcorr $\tilde Q$ 的格点计算转换为 PDF 有几种不同的途径。三个例子是：1) 先用式~(\ref{eq:fac-tq-orig0})傅里叶变换为准 PDF，再用式~(\ref{eq:fac-tq})；2) 先用式~(\ref{eq:pseudoRen})傅里叶变换为赝 PDF，再用式~(\ref{eq:ps-q-fact})；3) 先用式~(\ref{eq:io-Q-fact})匹配到坐标空间 PDF 的傅里叶变换 $Q(\zeta,\mu)$，再用式~(\ref{eq:FTq})的逆变换到 PDF。由于这些步骤的数值实现可能带有略微不同的系统误差，比较它们、或者使用多于一种方法以减小不确定性是有意义的。

根据第~\ref{sec:ope}~节的分析，要使欧氏分布的因子化公式成立，必须在小的距离 $z^2$ 与大的动量 $P^z$ 下计算同一\spcorr，以压低动力学与运动学的高扭度效应。对实际的格点计算而言，这意味着可用于提取 PDF 的 $(z, P^z)$ 数据点只有有限多个。

为说明这一点，考虑一个格距 $a=0.09$ fm 的 $48^3\times64$ 格点。空间关联的距离 $z$ 以 $a\sim1/2.2\ {\rm GeV}^{-1}$ 为单位，核子动量 $P^z$ 以 $2\pi/(48a)\sim0.29$ GeV 为单位。取 $\Lambda_{\rm QCD}\sim 0.3$ GeV。原则上靶质量修正可以被减除。若对整数 $m$ 与 $n$ 考虑各种值 $z=ma$ 与 $P^z=n*2\pi/(48a)$，那么为满足 $z\Lambda_{\rm QCD}\ll1$ 与 $P^z\gg \Lambda_{\rm QCD}$，必须有
\beq
 m\ll 11\,,\,\,\,\quad n\gg 1\,.
\eeq
要把高扭度修正控制在 $20\%$，即 $z^2\Lambda_{\rm QCD}^2\sim0.2,\ \Lambda_{\rm QCD}^2/P_z^2\sim0.2$，只能选取
\beq
m = \{0,1,2,3,4\}\,,\,\,\,\, n=\{3,4,5,6,\cdots\}\,,
\eeq
其中 $n$ 的最大值受当前格点模拟实际可行性限制。6 是文献~\cite{Orginos:2017kos}中达到的最大单位数。对准 PDF 计算，每个固定动量 $|P^z|$ 有 $4\times2+1=9$ 个有用数据点；对赝 PDF 计算，每个固定 $|z|$ 只有 $4\times2=8$ 个有用数据点。无论哪种情形，预计对 $z$ 或 $\zeta=zP^z$ 的直接傅里叶变换都会因坐标空间的截断而在 $x$ 空间产生振荡，并对小 $x$ 区域给出错误预言。这一现象最近在文献~\cite{Chen:2017mzz}的准 PDF 格点计算中已被观察到。近年来的工作已发展出消除截断效应振荡的方法~\cite{Lin:2017ani,Zhang:2017gau}（用于准 PDF），而大 $z$ 处的高扭度贡献仍需系统性修正。应当指出，以上只是高扭度修正的粗略估计，因为 $z^2\Lambda_{\rm QCD}^2$ 的前置因子可能小于 1。其实际重要性只能由格点模拟定量确定。

为充分利用所有有用数据点，可以根据微扰分析把它们演化到相同的 $z^2$ 或 $P^z$，这在文献~\cite{Orginos:2017kos,Karpie:2017bzm,Radyushkin:2017lvu}中针对\spcorr 已有研究。然而，由于\spcorr 在 $\ln z^2$ 或 $\ln P_z^2$ 中的演化方程遵循 $\zeta=zP^z$ 或 $z$ 的非定域卷积，做演化必须知道坐标空间的全部信息。在数据点有限的情况下，要么不确定度很大，要么不得不采用关于形状的模型依赖假设。

要提高任一方法的精度，唯一的出路是拥有更细的格距 $a$，从而获得更多满足 $|z|\ll \Lambda_{\rm QCD}^{-1}$ 的数据点和更大的核子动量 $P^z$。
随着 $P^z$ 增大，核子的价分布在 $z$ 方向上收缩，价夸克的空间关联收缩到更小的 $z$ 距离。若 $P^z$ 足够大，空间关联将在 $|z|< \Lambda_{\rm QCD}^{-1}$ 内迅速衰减，则来自傅里叶变换的截断误差将显著减小。另一方面，如果把该空间关联解释为\spcorr，其形状在洛伦兹助推下不会改变，因为它是 $\zeta=zP^z$ 与 $z^2$ 的标量函数。尽管如此，更细的格距 $a$ 允许在更宽的 $P^z$ 范围内进行计算，从而覆盖更大的 $\zeta=zP^z$ 值以减小截断误差。由于有用数据点的数目随 $1/a$ 二次增长，未来将出现系统误差可控的更精确格点计算。
